\let\textcircled=\pgftextcircled
\chapter{Impl\'ementation et r\'esultats}
\label{chap:implementation}

\textit{\initial{D}ans ce chapitre, on va présenter d'une part l'implémentation à travers les différentes étapes par lesquelles nous sommes passés; et, d'autre part, les résultats que l'on a obtenu à l'issu de ce stage.  }

\minitoc
\mtcselectlanguage{french}
\newpage

\section{Impl\'emetation}
\subsection{Cr\'eation d'une image avec supervisord}
Plus tard, il a aussi fallu g\'erer une autre pr\'eoccupation concernant l'image qui devait tourner \`a la fin. C'est-\`a-dire, soit l'API dans la base de donn\'ees, soit la base de donn\'ees dans l'API. Il a sembl\'e plus coh\'erent que ce soit l'API dans la base de donn\'ees car le serveur de base 
de donn\'ees doit tourner avant l'API pour que l'ensemble fonctionne correctement. D\`es lors, nous nous sommes plus int\'eress\'es en d\'etail \`a comment on pouvait avoir plusieurs processus dans un conteneur avec \cite{postgres-supervisord} et \cite{multiple-services-docker}.

Ensuite, nous avons commenc\'e \`a modifier notre Dockerfile pour que ce soit soit le conteneur serveur de base de donn\'ees qui re\c{c}oive l'API en copiant le dossier \lstinline{/app} g\'en\'er\'e lors de la compilation de notre API pour le coller dans un dossier du serveur de base de donn\'ees, mais \'egalement 
r\'earranger le Dockerfile de fa\c{c}on globale pour avoir une compilation effective. Connaissant la difficult\'e que ce serait de lancer les deux processus correctement en m\^eme temps, nous avons d\'ecid\'e de nous concentrer \`a d'abord faire marcher chacun d'eux correctement de fa\c{c}on individuelle. 
Ce qui f\^ut le cas de notre API suivant \cite{node-supervisor}.

Gr\^ace \`a \cite{node-supervisor}, on voit qu'il est n\'ecessaire d'avoir un serveur Nginx \footnote{que l'on installera et configurera avec un fichier de configuration} et un script o\`u on entrera les commandes n\'ecessaires pour le lancement de l'API. Cependant, sachant que notre API 
a \'et\'e compil\'ee avec \lstinline{WORKDIR /app} pour faciliter la r\'esolution des chemins relatifs et s'assurer d'\^etre dans le bon dossier pour effectuer des commandes. Ce que nous avons pu pallier en rempla\c{c}ant tous les chemins relatifs dans le code de l'API par des chemins absolus. 
De plus, on se retrouve \`a copier tout le dossier \lstinline{/app} donc aussi le dossier \lstinline{/node_modules} g\'en\'er\'e, ce qui peut entra\^iner des erreurs inattendues. Dans notre cas, il y a eu un probl\`eme avec le module \lstinline{dotenv} qui n'arrivait pas \`a se charger, ce qui 
nous a pouss\'e \`a commenter toute ligne o\`u \lstinline{dotenv} pouvait figurer tout en arr\^etant d'utiliser les variables d'environnement. En effectuant tout \c{c}a, jusqu'aux instructions de \cite{node-supervisor}, on arrive \`a lancer l'API.

Concernant la base de donn\'ees, nous avons d\'ecid\'e de comprendre le processus de d\'emarrage d'un serveur PostgreSQL, car ce sera \`a nous d'\'ecrire les commandes permettant de d\'emarrer le serveur. Pour ce faire, on utilise toujours le volume contenant les donn\'ees de notre base de donn\'ees charg\'ee. 
Cette approche g\'en\`ere beaucoup de conflits notamment avec le montage du volume et les probl\`emes de permissions associ\'es aux fichiers n\'ecessaires nous ont largement entrav\'e. 

C'est pourquoi, sachant que l'on manipule une image \lstinline{postgres}. Nous avons d\'ecid\'e de nous int\'eresser \`a la proc\'edure de lancement de cette derni\`ere en allant fouiller sur leur d\'ep\^ot GitHub \cite{postgres-docker-repo} o\`u nous avons d\'ecouvert que le lancement d'un conteneur postgres 
\'etait ex\'ecut\'e par le script \lstinline{docker-entrypoint.sh} qui prend comme argument \lstinline{postgres} tel que \lstinline[breaklines]{docker-entrypoint.sh postgres}. Et c'est cette commande que nous allons utiliser avec \lstinline{supervisor} pour lancer notre serveur PostgreSQL 
car il s'agit de la fa\c{c}on officielle de d\'emarrer.

Enfin, ayant effectu\'e toutes ces configurations, on compile notre image qu'on arrive \`a lancer avec la commande :
\begin{lstlisting}[frame=single, language=bash, breaklines]
docker run --name back -e PGDATA=postgres-data -p 3000:3000 
\end{lstlisting} 
car on a notre image qui continue \`a utiliser le volume contenant nos donn\'ees et exposant le port 3000 pour communiquer avec l'API. Lors de la premi\`ere initialisation du conteneur, le serveur PostgreSQL prend du temps pour se lancer et 
charger la base de donn\'ees qui entra\^ine l'\'echec de la connexion de l'API \`a la base de donn\'ees. Cependant, apr\`es avoir arr\^et\'e et relanc\'e le conteneur, tout fonctionne normalement.
\subsection{Notions de tests fonctionnels avec Hurl}
Pour ce faire, il sera question d'utiliser Hurl pour \'ecrire des requ\^etes HTTP \`a notre API. Mais pas seulement, on doit aussi \^etre en mesure de v\'erifier que la r\'eponse renvoy\'ee est correcte et m\^eme d'estimer le d\'elai de r\'eponse ou en imposer un. 
D'o\`u la notion de test, car si l'une des propri\'et\'es d\'efinies n'est pas v\'erifi\'ee, cela va automatiquement entrainer l'\'echec des tests.

Ainsi, nous avons d\'ecid\'e d'apprendre \`a utiliser Hurl, lou\'ee pour sa facilit\'e d'utilisation. En effet, on d\'efinit nos requ\^etes dans un format de texte tr\`es simple que l'on enregistre dans un fichier \`a l'extension \lstinline{.hurl} 
puis que l'on utilise avec la commande \lstinline!hurl test.hurl!

Si aucun fichier n'est fourni en entr\'ee, Hurl va consid\'erer l'entr\'ee standard tel qu'on ait : \lstinline[language=bash]{echo GET http://example.org | hurl}. De fa\c{c}on globale, Hurl se comporte comme cURL\footnote{Interface en ligne de commande, 
destin\'ee \`a r\'ecup\'erer le contenu d'une ressource accessible par un r\'eseau informatique.} et aura tendance \`a afficher le r\'esultat de la requ\^ete de telle sorte que l'on obtienne :
\begin{lstlisting}[language=bash, frame=single, title={Comportement cURL}, breaklines]
$ echo GET http://httpbin.org/get | hurl
{
"args": {},
"headers": {
"Accept": "*/*",
"Accept-Encoding": "gzip",
"Content-Length": "0",
"Host": "httpbin.org",
"User-Agent": "hurl/0.99.10",
"X-Amzn-Trace-Id": "Root=1-5eedf4c7-520814d64e2f9249ea44e0"
},
"origin": "1.2.3.4",
"url": "http://httpbin.org/get"
}
\end{lstlisting}

Sachant que notre objectif est de : d\'efinir des tests, pouvoir voir le nombre de tests effectu\'es, voir le nombre de tests pass\'es et \'evaluer le temps de r\'eponse sans avoir \`a afficher le contenu des requ\^etes. Il existe l'option \lstinline{--test} pour qu'on ait :
\lstinline{hurl --test *.hurl}.

Connaissant cela, il ne nous reste plus qu'\`a apprendre \`a \'ecrire des tests. Fort heureusement, Hurl met \`a notre disposition son format de fichier gr\^ace auquel on peut pleinement profiter de son potentiel. En effet, ce fichier consiste en une s\'erie de requ\^etes 
qui seront ex\'ecut\'ees tour \`a tour et dont on aura le r\'esultat.\\
\begin{lstlisting}[language=bash, frame=single, title={*.hurl}]
GET http://example.org/endpoint1
GET http://example.org/endpoint2
\end{lstlisting}

D\`es lors, on peut commencer \`a \'ecrire des tests en commen\c{c}ant par la requ\^ete puis le code de r\'eponse HTTP attendu.
\begin{lstlisting}[language=bash, frame=single, title={*.hurl}, breaklines]
# Avec cette requete, on s'attend a recevoir une reponse OK
GET http://example.org/endpoint1
HTTP 200
\end{lstlisting}

Bien que ce soit d\'ej\`a int\'eressant, notre objectif n'est pas seulement d'avoir le code de r\'eponse de la requ\^ete mais \'egalement de pouvoir v\'erifier son contenu. C'est pourquoi Hurl int\`egre \lstinline{asserts} comme l'illustre 
la figure~\ref{fig:xpath-asserts} dans le cadre d'un test avec du HTML utilisant \lstinline{XPath asserts}.
\begin{figure}[H]
\centering
\includegraphics[width=0.5\linewidth]{fig05/xpath.png}
\caption{Structure d'un assert pour le test du HTML \cite{Hurl}}
\label{fig:xpath-asserts}
\end{figure}

Par cons\'equent, on aurait pour un test qui v\'erifie si une page HTML renvoy\'ee par la requ\^ete a bien le titre "Morning".
\begin{lstlisting}[language=bash, frame=single, title={*.hurl}, breaklines]
GET http://example.org/endpoint

HTTP 200
[Asserts]
xpath "string(//head/title)" == "Morning"
\end{lstlisting}

On peut \'egalement v\'erifier l'ent\^ete d'une requ\^ete en utilisant \lstinline{header}. Cependant, l'API que l'on souhaite tester est REST et utilise par extension le format JSON, que l'on peut sp\'ecifier avec \lstinline{JSONPath assert}. Par exemple, 
on pourrait faire une requ\^ete \`a une API REST qui renvoie un objet repr\'esentant un utilisateur dont on aurait son nom et son pr\'enom.
\begin{lstlisting}[language=bash, frame=single, title={*.hurl}, breaklines]
GET http://example.org/endpoint

HTTP 200
[Asserts]
header "Content-Type" == "application/json; charset=utf8"
jsonpath "$.firstName" == "John"
jsonpath "$.lastName" == "Doe"     
\end{lstlisting}

De plus, on peut manipuler des tableaux d'objets en sp\'ecifiant l'index i correspondant \`a l'objet voulu dont on peut manipuler l'attribut via \lstinline{"$[i].attribut"}. On peut \'egalement d\'efinir les param\`etres de nos requ\^etes et 
un temps de r\'eponse. Par exemple, si l'on souhaite effectuer la requ\^ete \lstinline{http://example.org/endpoint/tabs?id=2} et v\'erifier qu'elle s'effectue en moins d'une seconde\footnote{Le temps est entr\'e en millisecondes.}, on peut \'ecrire : 
\begin{lstlisting}[language=bash, frame=single, title={*.hurl}, breaklines]
GET http://example.org/endpoint/tabs
[QueryStringParams]
id: 2

HTTP 200
[Asserts]
header "Content-Type" == "application/json; charset=utf8"
jsonpath "$.firstName" == "John"
jsonpath "$.lastName" == "Doe"     
duration < 1000
\end{lstlisting}

On dispose alors d'assez d'\'el\'ements pour nous consacrer \`a l'\'ecriture du g\'en\'erateur de charge dont le but est d'estimer la performance de notre conteneur sous une affluence variable de requ\^etes.

% Journ\'ee du 22/07/2024
% ~\newline\newline
\subsection{Cr\'eation d'une image sans volume Docker}
On doit faire en sorte que l'image Docker de la base de donn\'ees contienne directment les donn\'ees sans avoir recours \`a un volume. Pour ce faire, 
nous avons d\'ecid\'e de revoir nos fondamentaux notamment sur l'\'ecriture d'un Dockerfile ou ce que peut nous renseigner le d\'ep\^ot officiel de l'image Docker PostGIS/PostgreSQL utilis\'ee ~\cite{docker-postgres}.

Cela nous a amen\'e \`a nous rappeler du probl\`eme que l'on avait eu concernant la cr\'eation des utilisateurs avant le lancement du conteneur et le fait qu'il existait un dossier dans le syst\`eme de fichiers de
l'image Docker de PostgreSQL charg\'ee de l'initialisation\footnote{Lors du lancement du conteneur, ce dossier sera v\'erifi\'e et s'il y a des scripts, ils seront ex\'ecut\'es.} de scripts \lstinline{.sh} et 
\lstinline{.sql}. Il s'agit du dossier \lstinline{/docker-entrypoint-initdb.d}.

Cependant, on ne l'a pas directement utilis\'e, on a d'abord commenc\'e par lancer notre serveur de base de donn\'ees, ensuite on a eu \`a cr\'eer les utilisateurs n\'ecessaires et effectuer un 
\lstinline{docker commit}  de notre conteneur; afin d'obtenir une image de serveur ayant ces utilisateurs, pour faciliter l'initialisation de la base de donn\'ees. Cependant, 
au lancement de cette image, on a eu la r\'eponse selon laquelle les utilisateurs voulus n'existaient pas.
% \section{Journ\'ee du 23/07/2024}

\subsection{Tests de validit\'e et de performance}

Par la suite, nous avons commenc\'e \`a appliquer la m\'ethode pr\'econis\'ee pour l'\'elaboration du g\'en\'erateur de charge. On s'est alors mis \`a cr\'eer le fichier \lstinline{.hurl} n\'ecessaire 
pour \'evaluer la performance et valider la requ\^ete utilisant Hurl. Pour ce faire, on va utiliser la requ\^ete \url{http://<adresse-ip>:3000/geolocdpe/api/v0/dpe/get?topLeft=6692755.6320129
22,348070.5115176897&bottomRight=6686698.84754637,363102.02091083454} comme support. Gr\^ace \`a elle, on peut voir que notre API renvoie un tableau d'objets. Ainsi, pour v\'erifier la validit\'e, 
on n'a qu'\`a v\'erifier que tous les objets du tableau ont un id, ce que repr\'esente la ligne 8.

\begin{lstlisting}[language=bash, frame=single, title={load.hurl}, breaklines]
GET http://<adresse-ip>:3000/geolocdpe/api/v0/dpe/get
[QueryStringParams]
topLeft: 6692755.632012922,348070.5115176897
bottomRight: 6686698.84754637,363102.02091083454 

HTTP 200
[Asserts]
jsonpath "$.*.id" exists
\end{lstlisting}

Ayant \'ecrit ce fichier, on peut l'utiliser pour \'evaluer la performance de notre API en utilisant la commande :

\begin{lstlisting}[language=bash, frame=single, title={Commande Docker}, breaklines]
$ docker run --rm -v ./load.hurl:/tmp/load.hurl -w /tmp ghrc.io/orange-opensource/hurl:latest --test --color /tmp/load.hurl
\end{lstlisting}

On a donc le r\'esultat pr\'esent\'e \`a la figure~\ref{fig:result-req-hurl}. On voit qu'on a un temps de r\'eponse de 772 millisecondes. 
Cela repr\'esente la performance de notre API suivant une requ\^ete, sachant que la valeur peut changer si on la r\'eex\'ecute, 
on ne saurait la prendre comme valeur de r\'ef\'erence. C'est donc pour \c{c}a que l'on va adopter une approche na\"ive consistant \`a 
faire la moyenne de 10 requ\^etes et prendre la moyenne comme r\'ef\'erence. 
\begin{figure}[H]
\centering
\includegraphics[width=1\linewidth]{fig05/result_1_req_hurl.png}
\caption{R\'esultat d'une requ\^ete sur Hurl}
\label{fig:result-req-hurl}
\end{figure}

On va donc modifier notre fichier \lstinline{load.hurl} pour y mettre 10 requ\^etes et l'ex\'ecuter. On obtient le r\'esultat 
pr\'esent\'e \`a la figure~\ref{fig:result-10-req-hurl}. On voit qu'on a une dur\'ee de 7438 millisecondes. En divisant par 10, on a 743,8 millisecondes qui repr\'esente la performance d\'esir\'ee par requ\^ete et donc notre tol\'erance. 
\begin{figure}[H]
\centering
\includegraphics[width=1\linewidth]{fig05/result_10_req_hurl.png}
\caption{R\'esultat de 10 requ\^etes successives sur Hurl}
\label{fig:result-10-req-hurl}
\end{figure}

Ayant ainsi d\'efini notre tol\'erance, il est maintenant question de d\'efinir la charge maximale du syst\`eme utilisant hurl.

% \section{Journ\'ee du 24/07/2024}
% ~\newline\newline
Nous nous sommes ainsi document\'es sur l'utilisation de hurl et des outils dont il est inspir\'e : Siege \cite{siege-readme} et wrk \cite{wg-wrk}.

En effet, en lisant la documentation de hurl, nous l'avons trouv\'e insuffisante pour pleinement comprendre ce qu'elle proposait et en faire le tour. 
C'est dans cette optique nous avons d\'ecid\'e d'aller voir du c\^ot\'e de Siege et wrk, vu qu'il en est fait mention dans la documentation de hurl. 
On s'est rendu compte qu'ils fonctionnaient globalement de la m\^eme fa\c{c}on sauf que Siege nous permet de sp\'ecifier nos diff\'erentes requ\^etes 
dans un fichier texte tandis que wrk permet d'\'ecrire et d'ex\'ecuter des scripts Lua pour r\'ealiser les tests. 
Mais \'egalement, ils se distinguent au niveau des donn\'ees renvoy\'ees \`a la fin du test; bien qu'ils permettent de d\'efinir le nombre d'utilisateurs simul\'es, 
le nombre de threads, le nombre de requ\^etes et la dur\'ee. On va donc pr\'esenter ce que retourne chaque outil. On a notamment pour hurl.
    \begin{lstlisting}[frame=single, language=bash, title={Output de hurl},breaklines]
$ hurl "https://google.com" --calls=100 -p100 -f1000

Running 1 threads 100 parallel connections per thread with 100 requests per connection
+-----/-----+------+------+-----------+------+-----+-----+
|Cpted/Rqted|IdlKil|Errors|kBytesRecvd|Elpsed|Req/s|MB/s |
+-----/-----+------+------+-----------+------+-----+-----+
|  572/  665|     0|     0|     334.43| 1.00s|1118s|0.33s|
| 1000/ 1000|     0|     0|     257.73| 1.50s|668s |0.25s|
|
| RESULTS:             ALL
| fetches:             1000
| max parallel:        100
| bytes:               3.644410e+05
| seconds:             1.501000
| mean bytes/conn:     364.441000
| fetches/sec:         666.222518
| bytes/sec:           2.427988e+05
| HTTP response codes: 
| 200 -- 1000
    \end{lstlisting}

Puis pour wrk.
\begin{lstlisting}[frame=single, title={Output de wrk},breaklines]
$ wrk -t12 -c400 -d30s http://127.0.0.1:8080/index.html   

Running 30s test @ http://127.0.0.1:8080/index.html
12 threads and 400 connections
Thread Stats   Avg      Stdev     Max   +/- Stdev
Latency   635.91us    0.89ms  12.92ms   93.69%
Req/Sec    56.20k     8.07k   62.00k    86.54%
22464657 requests in 30.00s, 17.76GB read
Requests/sec: 748868.53
Transfer/sec:    606.33MB
\end{lstlisting}

Et enfin pour siege.

\begin{lstlisting}[frame=single,title={Output de siege},breaklines]
$ siege -u shemp.whoohoo.com/Admin.jsp -d1 -r10 -c25

..Siege 2.65 2006/05/11 23:42:16
..Preparing 25 concurrent users for battle.
The server is now under siege...done
Transactions: 250 hits
Elapsed time: 14.67 secs
Data transferred: 448000 bytes
Response time: 0.43 secs
Transaction rate: 17.04 trans/sec
Throughput: 30538.51 bytes/sec
Concurrency: 7.38
Status code 200: 250
Successful transactions: 250
Failed transactions: 0    
\end{lstlisting}

On peut ainsi \`a partir des r\'esultats retourn\'es choisir les caract\'eristiques que l'on souhaite \'etudier lors de nos tests.

On a effectu\'e une s\'erie de diff\'erentes requ\^etes de fa\c{c}on successive afin d'\'evaluer la performance. 
On s'est toujours retrouver dans une moyenne de 700 millisecondes. Cependant, en lan\c{c}ant deux processus diff\'erents 
simultan\'ement : l'un effectuant 10 requ\^etes pareilles et l'autre 10 requ\^etes diff\'erentes \`a notre API; 
on a vu que chaque processus a pris moins de 15 secondes.

% \section{Journ\'ee du 25/07/2024}
% ~\newline\newline
Nous avons continu\'e \`a utiliser hurl avec divers param\`etres pour tester notre image personnalis\'ee \lstinline{bd-api}, 
notamment le nombre de connexions parall\`eles envoyant un nombre donn\'e de requ\^etes. On \'ecrivait une commande dans 
le style de celle qui est pr\'esent\'ee en dessous o\`u on cr\'ee 10 connexions parall\`eles, chaque connexion effectuant 
la requ\^ete 2 fois durant un d\'elai de 30 secondes.

\begin{lstlisting}[frame=single, language=bash, title={Exemple de commande hurl}, breaklines]
$ hurl "http://<adresse-ip>:3000/geolocdpe/api/v0/dpe/get?topLeft=6692755.632012922,348070.5115176897&bottomRight=6686698.84754637,363102.02091083454" --calls=100 -p10 -f2 --timeout=30
\end{lstlisting}

D\`es lors, nous nous sommes rendus compte que l'on ne recevait plus de r\'eponse. On a donc d\'ecid\'e d'inspecter les journaux 
de notre conteneur en entrant \lstinline{docker logs bd-api} pour voir ce qu'il s'est pass\'e. On a vu qu'il y \'etait \'ecrit : 
\lstinline[breaklines]{error: could not resize shared memory segment "/PostgreSQL.222067968" to 4194304 bytes: No space left on device}.

Cela est une erreur provenant de notre base de donn\'ees PostgreSQL sp\'ecifiant que la m\'emoire partag\'ee est trop faible pour pouvoir r\'epondre 
\`a toutes les requ\^etes. En effet, l'extension utilis\'ee\footnote{PostGIS} utilis\'ee dans notre base de donn\'ees utilise intensivement la m\'emoire 
partag\'ee pour effectuer des calcus g\'eospatiaux ou stocker temporairement des r\'esultats interm\'ediaires. Ainsi, plusieurs requ\^etes en m\^eme temps 
ont vite satur\'e la m\'emoire. Ce probl\`eme peut \^etre palli\'e lors de la cr\'eation du conteneur comme le pr\'esente \cite{stack-not-resize-memory}; 
car Docker limite la taille partag\'ee d'un conteneur \`a 64 Mo, ce qu'on peut m\^eme augmenter jusqu'\`a 1 Go lors de la cr\'eation du conteneur. 
On va donc arr\^eter notre conteneur, le supprimer et le cr\'eer de nouveau en entrant la commande suivante :
\begin{lstlisting}[frame=single, language=bash, title={Cr\'eation de conteneur avec une plus grande m\'emoire partag\'ee},breaklines]
$ docker run --name bd-api -p 3000:3000 -e PGDATA=postgres-data -d --shm-size=1g bd-api
\end{lstlisting}

Par la suite, on a recommenc\'e \`a avoir des r\'eponses. Cependant, en augmentant le nombre de connexions parall\`eles jusqu'\`a 10 000, 
on a vu que hurl n'arrivait pas effectuer les requ\^etes en affichant le message d'erreur : 
\lstinline{Error threads[1]*parallelism[10000] > process fd resource limit[1024]}\footnote{Car il n'utilise qu'un seul thread 
ne pouvant contenir pas plus de 1024 connexions parall\`eles}. Parvenu jusqu'\`a ce niveau, il f\^ut question de se revoir et de red\'efinir ce que l'on voulait et comment y arriver.

% \section{Journ\'ee du 26/07/2024}
\subsection{Ecriture du g\'en\'erateur de charge}
Nous avons commenc\'e en nous rappelant de ce que l'on cherchait : combien le conteneur peut supporter de requ\^etes par seconde sachant qu'il doit r\'epondre en moins d'une seconde \`a 
chacune d'elle ? S'\'etant souvenu de \c{c}a, il \'etait maintenant question d'\'evaluer le temps moyen de r\'eponse suivant le nombre de requ\^etes par exemple. Pour ce faire, nous avons 
choisi une approche s\'equentielle o\`u on effectuait une requ\^ete suivant un d\'elai, d'abord chaque seconde puis chaque 900 ms jusqu'\`a 100 ms. Cela nous permet de quitter 
d'une requ\^ete par seconde \`a 10 requ\^etes par seconde. Pour ce faire, on va d'abord faire le cas pour une requ\^ete similaire puis s'inspirer de Hurl suivant \cite{concurrency-Hurl}. 
On va donc \'ecrire un script shell que l'on va nommer \lstinline{glob.sh} qui va nous permettre d'automatiser le processus :
\begin{lstlisting}[frame=single, language=bash, title={glob.sh}, breaklines]
#!/bin/bash
# glob.sh
#   1st argument: duration in seconds
#   2nd argument: name of the file to store the average time

run(){
# run
#   1st argument: time to sleep (delay)
#   2nd argument: duration
#   3rd argument: filename to store the average time

counter=0
sum=0
duration=$2
duration_int=$duration

while [ true ] 
do 
((counter = $counter + 1)) 
echo "Requete num: $counter" 
result=$(docker run --rm -v ./load.hurl:/tmp/load.hurl -w /tmp ghcr.io/orange-opensource/hurl:latest --test --color --json /tmp/load.hurl | jq .time)
((sum = $sum + $result)) 

sleep $1

duration=$(echo "$duration - $1" | bc)
duration_int=$(printf "%.0f" $duration)
echo "On a $duration avec $duration_int"
if [ $duration_int -lt 0 -o "${duration_int:0:1}" = "-" ]; then
    break
fi
done
echo -e "$(($sum / $counter))" >> $3
}

for(( i=0; i<10; i++ )); do
index=$(echo "1 - 0.$i" | bc)
run $index $1 $2
done

python3 glob.py $2
\end{lstlisting}

Pour \'ecrire ce script, on a d\^u installer \lstinline{jq}\footnote{Outil pour g\'erer le JSON en invite de commande} et \lstinline{bc}\footnote{Calculatrice} avec \lstinline{sudo apt install}. 
Ce script prend donc en entr\'ee la dur\'ee de notre test pour chaque tour et le nom du fichier o\`u doivent \^etre stock\'ees nos valeurs moyennes \`a chaque tour. Un tour ici, repr\'esente le d\'elai, 
par exemple une requ\^ete par seconde durant 30 secondes dont 30 requ\^etes et ainsi de suite pour le reste. Sanchant que l'on quitte d'une seconde \`a 100 millisecondes, on aura 10 tours. 
Enfin, \`a la fin de notre script, un programme Python sera appel\'e pour lire les valeurs dans le fichier qui lui est pass\'e en param\`etre et tracer une courbe repr\'esentant nos r\'esultats. 
Ce script Python que nous avons nomm\'e \lstinline{glob.py} est pr\'esent\'e tel que suit :
\begin{lstlisting}[language=python, frame=single, title={glob.py}, breaklines]
# glob.py 
#   argument: filename to take the data

import numpy as np 
import matplotlib.pyplot as plt 
import sys

def plot(filename):
rows = []
tab = []
with open(filename, 'r') as file:
rows = file.readlines()
for row in rows:
    tab.append(int(row))
x = np.array(list(range(1, 11)))
y = np.array(tab)

plt.plot(x, y, label='ms/req/s')

plt.title("Moyenne des requetes par seconde")
plt.xlabel("Req/s")
plt.ylabel("Temps moyen ms")
plt.grid(True)
plt.legend()

plt.savefig("result.png")

if __name__ == "__main__":
if len(sys.argv) > 1:
plot(sys.argv[1])
else:
print("Il manque le nom du fichier.")
\end{lstlisting}

Pour \'ecrire ce script, on a d\^u installer les paquets \lstinline{numpy}\footnote{Biblioth\`eque Python pour g\'erer les tableaux, matrices et calculs sur ces derniers.} et 
\lstinline{matplotlib}\footnote{Biblioth\`eque Python pour le trac\'e de courbes et graphes.} suivant la commande \lstinline{sudo apt install python3-<paquet>}. 
Il faudra d'abord s'assurer que notre script shell peut \^etre ex\'ecutable en entrant la commande \lstinline{chmod +x glob.sh}; puis, on peut lancer globalement 
sur une p\'eriode de 30 secondes, en entrant la commande suivante dans le dossier o\`u se trouvent nos deux scripts :
\begin{lstlisting}[language=bash, frame=single, title={Ex\'ecution de glob.sh}, breaklines]
$ sudo ./glob.sh 30 sample.txt
\end{lstlisting}

A la fin de l'ex\'ecution, les fichiers \lstinline{sample.txt} et \lstinline{result.png} vont se cr\'eer. On peut notamment voir la courbe sur la figure~\ref{fig:courbe_mesure_30s}.

\begin{figure}[H]
\centering
\includegraphics[width=1\linewidth]{fig05/result.png}
\caption{Courbe r\'esultante des mesures sur 30 s}
\label{fig:courbe_mesure_30s}
\end{figure}
De m\^eme, nos mesures sur chaque tour sont stock\'ees telles que :
\begin{lstlisting}[frame=single, title={sample.txt}]
783
729
748
739
740
741
736
744
739
737
\end{lstlisting}

% \section{Journ\'ee du 29/07/2024}
\subsection{Am\'elioration du g\'en\'erateur de charge avec g\'en\'eration de requ\^etes al\'eatoires}
Pour ce faire, nous avons analys\'e la structure de la requ\^ete faite \`a l'API. En effet, celle-ci utilise le syst\`eme de coordonn\'ees L93\footnote{Lambert 93, France m\'etropolitaine.}. 
Ce syst\`eme de coordonn\'ees repr\'esente les points en m sur deux dimensions : \texttt{x} et \texttt{y}. Notre API requiert sur le endpoint \lstinline{/get} deux points qui sont : 
le point haut gauche (\lstinline{topLeft}) et le point bas droite (\lstinline{bottomRight}) dont on sp\'eficie les coordonn\'ees s\'epar\'ees par des virgules telles que \lstinline{y,x}. 
Cela se repr\'esente dans la requ\^ete par \lstinline[breaklines]{http://10.144.208.247:3000/geolocdpe/api/v0/dpe/get?topLeft=y1,x1&bottomRight=y2,x2}.

Alors, il nous est venu l'id\'ee de nous donner un intervalle dans lequel nos coordonn\'ees seraient g\'en\'er\'ees. Sachant que notre base de donn\'ees contient des donn\'ees relatives \`a la France, 
il f\^ut donc question de choisir un intervalle s\^ur pour lequel il n'y aurait pas d'erreur. Pour ce faire, nous nous sommes aid\'es de \cite{geoportail} qui nous a permis de choisir nos points 
comme l'illustre la figure~\ref{fig:limites_gen_alea}.
\begin{figure}[H]
\centering
\includegraphics[width=0.5\linewidth]{fig05/limite.png}
\caption{Limites pour la g\'en\'eration al\'eatoire}
\label{fig:limites_gen_alea}
\end{figure}

Bien que cet intervalle ne recouvre pas toute la France, il nous permet de tester l'API et de toujours avoir des requ\^etes qui marchent. Ayant choisi nos points, on note les coordonn\'ees :
\begin{itemize}
\item [\textbullet] \texttt{topLeft (x = 441423.15 et y = 6763025.44)}
\item [\textbullet] \texttt{bottomRight (x = 500455.83 et y = 6638807.60)}
\end{itemize}

Avec cette approche, il s'agit de modifier notre code et ajouter une fonction qui va g\'en\'erer al\'eatoirement un nombre dans les intervalles d\'efinis puis \'ecrire dans le fichier 
utilis\'e pour effectuer des requ\^etes\footnote{Le fichier Hurl load.hurl.}. On va donc ajouter les lignes qui suivent dans notre script :
\begin{lstlisting}[frame=single, language=bash, title={Lignes \`a ajouter dans glob.sh}, breaklines]
write_file() {
text="GET http://10.144.208.247:3000/geolocdpe/api/v0/dpe/get\n[QueryStringParams]\ntopLeft: $2,$1\nbottomRight: $4,$3\n\nHTTP 200\n[Asserts]\njsonpath \"\$.*.id\" exists"
echo -e "$text" > "load.hurl"
}

x1=$(python3 generateCoordinates.py 1)
x2=$(python3 generateCoordinates.py 1)
y1=$(python3 generateCoordinates.py 2)
y2=$(python3 generateCoordinates.py 2)

write_file $x1 $y1 $x2 $y2
\end{lstlisting}

Ces lignes repr\'esentent une fonction qui prend en param\`etre les 04 valeurs de coordonn\'ees et \'ecrit dans le fichier \texttt{load.hurl} dont est issue la requ\^ete. Ceci \'etant dit, 
on a recours \`a Python pour g\'en\'erer al\'eatoirement la coordonn\'ee via le module \texttt{random} dans le script \texttt{generateCoordinates.py}.
\begin{lstlisting}[frame=single, language=python, title={generateCoordinates.py}, breaklines]
import random
import sys
def generate_x():
num = random.uniform(441423.15, 500455.83)
print(num)
def generate_y():
num = random.uniform(6638807.60, 6763025.44)
print(num)
if __name__ == "__main__" :
if len(sys.argv) > 1:
if sys.argv[1] == "1":
    generate_x()
elif sys.argv[1] == "2":
    generate_y()
else:
print("Pas d'argument.")
\end{lstlisting}

On passe en argument de notre script Python soit 1 ou 2 selon que l'on veuille g\'en\'erer x ou y. Puis viens l'heure du test, on va donc relancer notre script avec les diff\'erentes 
modifications et avoir la courbe pr\'esent\'ee \`a la figure~\ref{fig:courbe_aleatoire_10s}.
\begin{figure}[H]
\centering
\includegraphics[width=1\linewidth]{fig05/courbe_aleatoire_10s.png}
\caption{Courbe avec des requ\^etes al\'eatoires sur 10 secondes}
\label{fig:courbe_aleatoire_10s}
\end{figure}

On voit que peu importe le d\'elai, chacune de nos requ\^etes est ex\'ecut\'ee en moins d'une seconde, qu'\`a 0.1 seconde\footnote{C'est-\`a-dire 10 requ\^etes pas seconde.} 
on se retrouve \`a 152 ms en moyenne pour 100 requ\^etes effectu\'ees durant cette p\'eriode. Cependant, l'on se souvient que l'intervalle choisit est petit et 
qu'il ne couvre pas toute la France, c'est pourquoi il sera question prochainement de l'agrandir.

% \section{Journ\'ee du 30/07/2024}
\subsection{Am\'elioration du g\'en\'erateur en augmentant la surface couverte}
Nous avons continu\'e \`a \'evaluer notre image Docker \texttt{bd-api} en modifiant notre script \texttt{glob.sh}. En effet, on a pris qu'une petite portion de la France 
afin de tester comment notre conteneur r\'eagissait avec diff\'erentes requ\^etes suivant des d\'elais vari\'es. Ainsi, on a commenc\'e par \'etendre notre intervalle comme le pr\'esente la figure~\ref{fig:perimetre_etendu_France}.
\begin{figure}[H]
\centering
\includegraphics[width=1\linewidth]{fig05/perimetre_etendu_France.png}
\caption{Intervalle \'etendu \`a plus grande portion de la France}
\label{fig:perimetre_etendu_France}
\end{figure}

On peut y voir notre \texttt{topLeft} global et notre \texttt{bottomRight} global qui vont repr\'esenter l'intervalle dans lequel nos points de coordonn\'ees seront g\'en\'er\'es. Cependant, lors de la g\'en\'eration, 
il peut arriver que \texttt{bottomRight} soit en haut \`a gauche et \texttt{topLeft} soit en bas \`a droite, ce qui n'est pas un comportement voulu. Pour y rem\'edier, on va apporter de l\'eg\`eres modifications 
\`a notre script Python \texttt{generateCoordinates.py}.
\begin{lstlisting}[language=python, frame=single, title={generateCoordinates.py}, breaklines]
def generate_x():
num = random.uniform(454711.26, 1017778.19)
num2 = random.uniform(num, 1017778.19)
print(num, " ", num2)

def generate_y():
num = random.uniform(6326985.52, 6904672.61)
num2 = random.uniform(6326985.52, num)
print(num, " ", num2)  
\end{lstlisting}

En effet, on s'est rendu compte que plus on se d\'epla\c{c}ait vers la droite que la valeur de x augmentait tandis que la valeur de y diminuait lorsqu'on se d\'epla\c{c}ait vers le bas. C'est la raison pour laquelle, 
on modifie les limites de notre intervalle pour faire en sorte qu'ayant pris un x ou un y donn\'e, que la valeur de l'autre coordonn\'ee soit acceptable.

Apr\`es avoir effectu\'e ces modifications, nous avons relanc\'e notre script en sp\'ecifiant comme dur\'ee 30 secondes. On a obtenu les r\'esultats pr\'esent\'es \`a la figure~\ref{fig:req_30s_aleat_perimetre_etendu}.
\begin{figure}[H]
\centering
\includegraphics[width=1\linewidth]{fig05/req_30s_aleat_perimetre_etendu.png}
\caption{R\'esultats sur 30 secondes avec requ\^etes g\'en\'er\'ees al\'eatoirement}
\label{fig:req_30s_aleat_perimetre_etendu}
\end{figure}

On peut y voir que notre conteneur prend en moyenne plus d'une seconde pour r\'epondre \`a une requ\^ete selon le d\'elai, sauf pour chaque 0.8 seconde\footnote{Environ 1 requ\^ete par seconde.} o\`u se retrouve \`a 950 ms.
% \section{Journ\'ee du 31/07/2024}
En analysant les r\'esultats obtenus, on s'est interrog\'e sur leur objectivit\'e. En effet, les requ\^etes sont g\'en\'er\'ees de fa\c{c}on al\'eatoire; bien qu'on arrive \`a avoir de bons \texttt{topLeft} et \texttt{bottomRight}, 
on peut se retrouver \`a cr\'eer un rectangle trop grand. Parce que l'objectif premier est de retrouver un bien associ\'e \`a une annonce immobili\`ere, qui pourrait \^etre d'une surface bien moindre.

% \section{Journ\'ee du 02/08/2024}
\subsection{Utilisation de requ\^etes r\'eelles}
Toujours dans l'id\'ee d'estimer au mieux la charge du syst\`eme via un comportement r\'eel, on a laiss\'e tomber l'approche utilisant les requ\^etes g\'en\'er\'ees al\'eatoirement pour \`a la place choisir 30 requ\^etes r\'eelles 
que nous allons faire tourner al\'eatoirement. Pour ce faire, on a d\^u avoir recours \`a l'interface graphique accessible \`a \texttt{http://10.144.208.198:4200/} comme l'illustre la figure~\ref{fig:interface}. 
\begin{figure}[H]
\centering
\includegraphics[width=1\linewidth]{fig05/interface.png}
\caption{Interface de l'application}
\label{fig:interface}
\end{figure}

On peut voir sur la figure~\ref{fig:interface}, la console o\`u on peut copier la requ\^ete effectu\'ee lorque l'on se d\'eplace sur la carte. En effet, la requ\^ete est effectu\'ee suivant les dimensions de la carte qui est mat\'erialis\'ee; 
donc, pour g\'en\'erer une requ\^ete, on a qu'\`a se d\'eplacer et copier ce que l'on a en console.

La requ\^ete effectu\'ee par l'interface graphique \`a partir de la carte renvoie 500 lieux de fa\c{c}on al\'eatoire. Cela veut dire que plus la surface est grande et plus on risque de ne plus se retrouver avec les m\^emes points; 
cela entraine donc un temps de r\'eponse variant car il n'y a pas vraiment de mise en cache. Dans cette optique, on a d\'ecid\'e de choisir des requ\^etes qui renvoient au plus 500 lieux en zoomant la carte au maximum et 
de modifier notre script pour l'ex\'ecuter chaque 100 millisecondes mais modifier \`a chaque fois l'intervalle de la surface souhait\'ee. Pour ce faire, on va modifier la requ\^ete faite \`a l'API pour utiliser l'endpoint 
\lstinline{/dpe_with_par} qui sera \'ecrite dynamiquement dans le fichier \texttt{load.hurl}.

\begin{lstlisting}[language=bash, frame=single, title={load.hurl}, breaklines]
GET http://10.144.208.247:3000/geolocdpe/api/v0/dpe/dpe_with_par
[QueryStringParams]
topLeft: 6692949.509452589,854553.4328283346
bottomRight: 6692812.060103242,854722.7890251336
surface: 30100,30200

HTTP 200
[Asserts]
jsonpath "$.*.id" exists
jsonpath "$.*" count <= 500
\end{lstlisting}

On va modifier notre script que l'on va se retrouver \`a le lancer en entrant la commande : 
\begin{lstlisting}[frame=single]
$ sudo ./glob.sh 30 values.txt 100
\end{lstlisting}

O\`u comme premier param\`etre, on a la dur\'ee souhait\'ee en secondes, puis en second param\`etre le fichier o\`u stocker le temps moyen de r\'eponse de chaque requ\^ete et enfin l'intervalle choisi. 
Parce que l'on va d'abord commencer avec un minimum de 100 puis augmenter de la valeur de l'intervalle \`a chaque 100 millisecondes. Donc, pour les premi\`eres requ\^etes, on aura une surface 
comprise entre 100 et 200 puis 200 et 300; et ainsi de suite.

De plus, on a dans un premier temps pris dans 15 villes fran\c{c}aises 2 zones de coordonn\'ees :
\begin{itemize}
\item [\textbullet] Paris 
\item [\textbullet] Marseille 
\item [\textbullet] Lyon 
\item [\textbullet] Toulouse 
\item [\textbullet] Nice 
\item [\textbullet] Bordeaux 
\item [\textbullet] Lille 
\item [\textbullet] Nantes 
\item [\textbullet] Strasbourg 
\item [\textbullet] Montpellier 
\item [\textbullet] Rennes 
\item [\textbullet] Grenoble 
\item [\textbullet] Angers 
\item [\textbullet] Tours 
\item [\textbullet] Dijon
\end{itemize}

Nous avons ainsi enregistr\'e les diff\'erentes coordonn\'ees dans un fichier texte que nous avons nomm\'e \texttt{req.txt}. Apr\`es ex\'ecution de ce script, 
on se retrouve la courbe que pr\'esente la figure~\ref{fig:req_15_villes_France_10s}.
\begin{figure}[H]
\centering
\includegraphics[width=1\linewidth]{fig05/req_15_villes_France_10s.png}
\caption{Courbe des 30 requ\^etes sur 10 secondes}
\label{fig:req_15_villes_France_10s}
\end{figure}

% \section{Journ\'ee du 05/08/2024}
% ~\newline\newline
Nous avons continu\'e d'\'evaluer notre syst\`eme en modifiant l'intervalle de la surface de 100 \`a 5 $m^2$. Cependant, 
face \`a la rapidit\'e de la r\'eponse, nous nous sommes demand\'es combien de lieux \'etaient effectivement renvoy\'es sachant qu'il 
y avait une condition sur le nombre de lieux. Pour ce faire, nous avons d\'ecider de modifier notre fichier \texttt{load.hurl} 
afin de capturer\footnote{Stocker dans une variable, langage propre \`a Hurl.} le nombre de lieux retourn\'es via \cite{capture-Hurl}.
\begin{lstlisting}[frame=single, title={load.hurl}, breaklines]
GET http://10.144.208.247:3000/geolocdpe/api/v0/dpe/dpe_with_par
[QueryStringParams]
topLeft: 6692949.509452589,854553.4328283346
bottomRight: 6692812.060103242,854722.7890251336
surface: 150,151

HTTP 200
[Captures]
locations: jsonpath "$.*" count
[Asserts]
jsonpath "$.*.id" exists
jsonpath "$.*" count <= 500
\end{lstlisting}

Ainsi, pour voir le nombre de lieux retourn\'es, on entre la commande suivante dans le terminal :
\begin{lstlisting}[frame=single, language=bash, breaklines]
$ sudo docker run --rm -v ./load.hurl:/tmp/load.hurl -w /tmp ghcr.io/orange-opensource/hurl:latest --test --color --json /tmp/load.hurl | jq . entries[0].captures[0].value
\end{lstlisting}

La r\'eponse nous est pr\'esent\'ee sur la figure~\ref{fig:nb_locations_returned}. On voit que le nombre de lieux est \'egal \`a 0.
\begin{figure}[H]
\centering
\includegraphics[width=1\linewidth]{fig05/nb_locations_returned.png}
\caption{R\'eponse concernant le nombre de lieux retourn\'es}
\label{fig:nb_locations_returned}
\end{figure}

% \section{Journ\'ee du 06/08/2024}
\subsection{Mise en \'evidence du cache de PostgreSQL}
Nous avons un peu mieux cern\'e ce que l'on voulait tester. Au contact de l'interface graphique, nous nous sommes rendus compte que chaque mouvement de 
la carte entra\^inait une nouvelle requ\^ete que l'on pouvait voir dans la console. Ainsi, on comprend mieux pourquoi \'evaluer le nombre de requ\^etes qu'un 
conteneur peut g\'erer en moins d'une seconde. Cela est tr\`es int\'eressant dans la mesure o\`u rien qu'en utilisant l'interface graphique que l'utilisateur 
effectuait une succession de requ\^etes qui va se mat\'erialiser par des points repr\'esent\'es sur la carte. Si une requ\^ete prend trop de temps, 
cela peut entraver l'exp\'erience utilisateur.

Suivant cette logique, nous nous sommes int\'eress\'es \`a comment PostgreSQL mettait en place la mise en cache. Cela nous a permit de nous rendre compte que PostgreSQL 
n'avait pas de cache int\'egr\'e mais utilisait plut\^ot la m\'emoire partag\'ee du syst\`eme. D\`es lors, il f\^ut question de tester son impact, afin de v\'erifier si cette 
mise en cache influait sur le temps de r\'eponse des requ\^etes.

Pour v\'erifier cela, nous avons d'abord commenc\'e par voir comment on pouvait vider le cache de notre conteneur. On a vu que l'on pouvait entrer la succession de commande :
\begin{lstlisting}[frame=single, language=bash, breaklines]
$ docker exec -it bd-api bash
id-container> service stop postgresql
id-container> sync
id-container> echo 3 > /proc/sys/vm/drop_caches
id-container> service start postgresql
\end{lstlisting}

Cela s'est sold\'e par une erreur car on a d\'ecouvert que le cache utilis\'e par un conteneur est celui du syst\`eme, donc qu'il fallait ex\'ecuter la commande suivance au niveau du syst\`eme :
\begin{lstlisting}[frame=single, language=bash, breaklines]
$ sudo sh -c "/usr/bin/echo 3 > /proc/sys/vm/drop_caches"    
\end{lstlisting}

% \section{Journ\'ee du 07/08/2024}
% ~\newline\newline
On a continu\'e \`a \'evaluer \`a quel point la mise en cache optimisait le temps de r\'eponse. Pour ce faire, nous avons d'abord commenc\'e par vider le cache, 
puis nous avons ex\'ecut\'e le programme une premi\`ere fois puis une seconde fois pour voir la diff\'erence. Et enfin, on a ex\'ecut\'e le programme une troisi\`eme fois en vidant le cache \`a chaque requ\^ete. 
Cela nous a donn\'e les valeurs en millisecondes pr\'esent\'ees sur la figure~\ref{fig:val1_val2_val3}.
\begin{figure}[H]
\centering
\includegraphics[width=1\linewidth]{fig05/val1_val2_val3.png}
\caption{V\'erification de la mise en cache sur 30 requ\^etes distinctes}
\label{fig:val1_val2_val3}
\end{figure}

On voit que cela ne nous permet pas v\'eritablement de conclure, c'est pourquoi, nous avons plut\^ot d\'ecid\'e d'ex\'ecuter une requ\^ete avec un temps de r\'eponse significatif. 
On a choisi la requ\^ete \lstinline[breaklines]{http://10.144.208.247:3000/geolocdpe/api/v0/dpe/get?topLeft=6692755.632012922,348070.5115176897&bottomRight=6686698.84754637,363102.02091083454} 
qui r\'epond en moyenne en 700 millisecondes. Nous avons donc pris cette requ\^ete que nous avons ex\'ecut\'e 30 fois tout en vidant le cache, le r\'esultat est illustr\'e sur la figure~\ref{fig:same_req_without_cache}.
\begin{figure}[H]
\centering
\includegraphics[width=1\linewidth]{fig05/same_req_without_cache.png}
\caption{Requ\^ete significative effectu\'ee 30 fois sans cache}
\label{fig:same_req_without_cache}
\end{figure}

Par la suite, nous avons d\'ecid\'e de voir quel r\'esultat on aurait avec la mise en cache; nous avons donc comment\'e la commande pour mettre en cache et r\'eex\'ecuter 
le programme pour avoir ce qui est pr\'esent\'e sur la figure~\ref{fig:same_req_30x_with_cache}.
\begin{figure}[H]
\centering
\includegraphics[width=1\linewidth]{fig05/same_req_30x_with_cache.png}
\caption{Requ\^ete significative effectu\'ee 30 fois avec cache}
\label{fig:same_req_30x_with_cache}
\end{figure}

On voit clairement l'effet de la mise en cache sur notre syst\`eme. Ayant mis en exergue l'impact de la mise en cache, 
nous avons d\'ecid\'e d'ex\'ecuter notre programme 10 secondes en prenant 30 points dans 15 villes fran\c{c}aises, le r\'esultat est illustr\'e par la figure~\ref{fig:req_15_villes_fran\c{c}aises_10s_sans_cache}.
\begin{figure}[H]
\centering
\includegraphics[width=1\linewidth]{fig05/req_15_villes_francaises_10s_sans_cache.png}
\caption{Requ\^etes sur 15 villes fran\c{c}aises durant 10 secondes sans cache}
\label{fig:req_15_villes_fran\c{c}aises_10s_sans_cache}
\end{figure}

% \section{Journ\'ee du 08/08/2024}
% ~\newline\newline
Nous avons continu\'e \`a mettre en exergue la mise en cache. Pour ce faire, nous avons d\'ecid\'e de vider le cache apr\`es 5 requ\^etes. On a d\^u modifier la boucle principale de notre script \texttt{glob.sh}.
\begin{lstlisting}[frame=single, language=bash, breaklines]
req=0
for(( i=0; i<30; i++ )); do
((req = $req + 1))
echo "C'est $req"
run 1 $1 $2 $i
if [ $req -eq 5 ];then
sh -c "/usr/bin/echo 3 > /proc/sys/vm/drop_caches"
echo "Cache vide."
req=0
fi
done
\end{lstlisting}

Par la suite, nous avons d'abord lanc\'e notre script sur les requ\^etes sur 15 villes fran\c{c}aises, ce qui est repr\'esent\'e sur la figure~\ref{fig:cache_vide_chaque_5s}.
\begin{figure}[H]
\centering
\includegraphics[width=1\linewidth]{fig05/cache_vide_chaque_5s.png}
\caption{Requ\^etes sur 15 villes fran\c{c}aises avec cache vid\'e apr\`es 5 requ\^etes}
\label{fig:cache_vide_chaque_5s}
\end{figure}

Puis, nous avons lanc\'e un script effectuant une requ\^ete significative apr\`es 5 secondes, on se retrouve avec le r\'esultat sur la figure~\ref{fig:cache_vide_req_significative_chaque_5s}.
\begin{figure}[H]
\centering
\includegraphics[width=1\linewidth]{fig05/cache_vide_req_significative_chaque_5s.png}
\caption{Cache vid\'e apr\`es 5 requ\^etes sur une requ\^ete significative}
\label{fig:cache_vide_req_significative_chaque_5s}
\end{figure}

On peut clairement voir l'impact de la mise en cache car avec les requ\^etes 1, 6, 11, 16, 21 et 26 qui se passent apr\`es avoir vid\'e le cache.

% \section{Journ\'ee du 09/08/2024}
% ~\newline\newline
\subsection{Ajout de la table parcelles \`a l'image Docker finale}
Nous nous sommes consacr\'es \`a la cr\'eation d'une image comportant l'API et uniquement la table \texttt{parcelles}. 
Pour pouvoir mettre cela en place, il aura d'abord \'et\'e n\'ecessaire de r\'ealiser un nouveau dump pour cette table sp\'ecifiquement.

Par la suite, on a recours \`a l'image \texttt{postgis/postgis} pour cr\'eer un conteneur o\`u on aura copi\'e notre dump SQL. Puis on a cr\'e\'e les utilisateurs \texttt{cestou} et \texttt{express} 
avec \lstinline[language=sql, breaklines]{CREATE USER}. Et enfin, on aura charg\'e cette table en effectuant la commande : 
\begin{lstlisting}[frame=single, language=bash, breaklines]
$ docker exec -it first-parcelles psql -U j2m -f /var/lib/postgresql/parcelles_new_dump.sql
\end{lstlisting}

Une fois que cette table aura \'et\'e charg\'ee, on va stocker les donn\'ees dans un volume avant de faire un commit afin de cr\'eer une image personnalis\'ee :
\begin{lstlisting}[frame=single, language=bash, breaklines]
$ docker exec -it first-parcelles -r /var/lib/postgresql/parcelles_new_dump.sql parcelles data
$ docker commit first-parcelles bd-parcelles:complete
\end{lstlisting}

D\`es lors, on va modifier le code de l'API pour mettre le port \`a 3001 dans \texttt{nginx.conf} et \texttt{index.ts} puis on a comment\'e \texttt{dpe.route.ts} pour avoir :
\begin{lstlisting}[frame=single, breaklines]
// router.get('/get', method.getPointsInZone)
router.get('/parcelles',method.getParcelles)
// router.get('/dpe_with_par', method.getDpeWithParcelles)
\end{lstlisting}

Ensuite, on a modifi\'e le Dockerfile qu'on avait initialement cr\'e\'e pour mettre l'image \lstinline{bd-parcelles:complete} et on a compil\'e l'image. 
Puis, on a cr\'e\'e le conteneur avec :
\begin{lstlisting}[frame=single, language=bash, breaklines]
$ docker run -d --name container-api-parcelles -e PGDATA=parcelles-data -p 3001:3001 api-parcelles
\end{lstlisting}

Et pour tester le fonctionnement effectif de notre conteneur, on va faire une requ\^ete sur le endpoint \texttt{/parcelles} en entrant 
l'URL \url{http://10.144.208.247:3001/geolocdpe/api/v0/dpe/parcelles?topLeft=6692755.632012922,348070.5115176897&bottomRight=6686698.84754637,363102.02091083454&surface=10,100}. 
Et on aura le r\'esultat pr\'esent\'e par la figure~\ref{fig:test_req_parcelles}.
\begin{figure}[H]
\centering
\includegraphics[width=1\linewidth]{fig05/test_req_parcelles.png}
\caption{Test du endpoint parcelles}
\label{fig:test_req_parcelles}
\end{figure}

\section{R\'esultats}
    \subsection{Images Docker}
    \begin{figure}[H]
        \centering
        \includegraphics[width=1\linewidth]{fig05/Capture d'ecran 2024-07-16 051223.png}
        \caption{R\'ecapitulatif des diff\'erentes images Docker}
        \label{fig:recap-images}
    \end{figure}
On peut y voir :
    \begin{itemize}
        \item [\textbullet] \textbf{bd-api} : correspondant \`a notre image Docker personnalis\'ee permettant le lancement des services du serveur PostgreSQL et de l'API. C'est l'image finale obtenue en fin de semaine.
        \item [\textbullet] \textbf{postgres-dpe} : correspondant au commit de la base de donn\'ees charg\'ee avec le dump SQL.
        \item [\textbullet] \textbf{api-image} : correspondant \`a l'image Docker de l'API Express. 
        \item [\textbullet] \textbf{first} : image Docker du serveur PostgreSQL ayant le dump SQL import\'e. 
        \item [\textbullet] \textbf{registry} : image Docker permettant la cr\'eation d'un registre local \`a l'instar de Docker hub, dans le but d'avoir un emplacement de stockage local pour nos images afin de pouvoir y acc\'eder \`a distance.\\
    \end{itemize}
    \subsection{Comparaison des images}
    Nous nous sommes lanc\'es dans l'\'etude comparative de nos images Docker : la premi\`ere, r\'esultant d'un \lstinline{docker commit} utilisant un volume que nous avons nomm\'e \lstinline{postgres-dpe}; et la seconde, qui int\`egre directement le dump SQL, nomm\'ee \lstinline{second}. Pour la lisibilit\'e, on a choisi de pr\'esenter nos r\'esultats sous forme de tableau que pr\'esente le tableau~\ref{tab:tab1}.


    \begin{table}[H]
        \centering
        \begin{tabular}{c|c|c}
    \textbf{Caract\'eristiques} & \textbf{postgres-dpe} & \textbf{second} \\
    \hline
    Taille image & 15,4 Go & 6, 59 Go \\
    Temps d'initialisation & 2min33s & 28s \\
    Temps de lancement & 0s & 0s \\
    Temps d'arr\^et & 0s & 0s \\
    Temps de destruction & 0s  & 0s \\
    Persistence des donn\'ees & Oui & Non \\
    Utilisation concurrente & Oui & Oui \\
        \end{tabular}
        \caption{Tableau comparatif des images postgres-dpe et second}
        \label{tab:tab1}
    \end{table}
    
    Maintenant, il nous revient de sp\'ecifier quelques concepts :
    \begin{itemize}
        \item [\textbullet] \textbf{Temps d'initialisation :} temps n\'ecessaire lors du premier lancement du conteneur apr\`es \lstinline{docker run} pour charger les donn\'ees et rendre la base de donn\'ees accessible.
        \item [\textbullet] \textbf{Temps de lancement :} temps n\'ecessaire pour lancer le conteneur apr\`es l'avoir pr\'ealablement cr\'e\'e avec \lstinline{docker start}.
        \item [\textbullet] \textbf{Temps d'arr\^et :} temps que met le conteneur \`a s'arr\^eter apr\`es \lstinline{docker stop}.
        \item [\textbullet] \textbf{Temps de destruction :} temps que met le conteneur \`a \^etre d\'etruit arp\`es la commande \lstinline{docker rm}.
        \item [\textbullet] \textbf{Utilisation concurrente :} qui renvoie \`a la capacit\'e de cr\'eer plusieurs conteneurs utilisant la m\^eme image en  m\^eme temps.\\
    \end{itemize}
    
    Nous tenons aussi \`a ajouter que le temps de 0 seconde fait r\'ef\'erence au fait que l'op\'eration a pris moins d'une seconde. \\
    \subsection{D\'eploiement Kubernetes}
A l'issue des diff\'erents tests r\'ealis\'es, on a pu constater que le syst\`eme g\'erait plut\^ot bien sans mise en cache. Ce qui nous a permis de le d\'eployer sereinement via Kubernetes via 
l'usage de registre priv\'e s\'ecuris\'e sur le serveur qui a \'et\'e mis \`a notre disposition. En effet, on l'a d\'eploy\'e avec un nombre de 3 replicas sur un service de type \texttt{LoabBalancer} afin 
que la charge soit \'equitablement r\'epartie.
    \subsection{Production d'un tutoriel}
Dans le but d'encapsuler les \'etapes de conteneurisation avec Docker et du d\'eploiement avec Kubernetes, on a cr\'e\'e un tutoriel au format texte qui r\'esume la proc\'edure en passant par la cr\'eation du registre jusqu'au d\'eploiement.